\documentclass[10pt,a4paper]{amsart}
\usepackage[margin=19mm]{geometry}
\usepackage{amsmath,amssymb,mathtools}
\usepackage[scr=rsfs,cal=euler]{mathalfa}
\usepackage{xfrac}
\usepackage[unicode,hidelinks,bookmarks=false]{hyperref}
\newcommand{\ie}{i.e.\ }
\newcommand{\cf}{cf.\ }
\newcommand{\resp}{respectively\ }
\newcommand{\Subsection}{\subsection}
\providecommand{\nobreakdash}{\nobreak\hbox{-}}
\setlength{\emergencystretch}{2em}
\multlinegap0pt
\usepackage{bm}
\usepackage{bbm}
\usepackage{dsfont}
\usepackage{xspace}
\usepackage{cancel}
\usepackage{mathtools}
\usepackage[shortlabels]{enumitem}
\usepackage{amsthm}
\makeatletter
\@ifundefined{theorem}{\newtheorem{theorem}{Theorem}}{}
\@ifundefined{definition}{\newtheorem{definition}[theorem]{Definition}}{}
\@ifundefined{proposition}{\newtheorem{proposition}[theorem]{Proposition}}{}
\makeatother
\newcommand{\Z}{\mathbb{Z}}
\title[A tree bijection for rigid quadrangulations]{A tree bijection for rigid quadrangulations}
\author{Bart Zonneveld}
\date{Source: arXiv:2509.24787v1 (CC BY 4.0)}
\begin{document}
\maketitle

\noindent\emph{Source and licence.} A tree bijection for rigid quadrangulations. Version arXiv:2509.24787v1 (CC BY 4.0), distributed under the Creative Commons Attribution 4.0 International licence, as recorded in the arXiv licence metadata for this exact version and shown on the abstract page https://arxiv.org/abs/2509.24787v1. Full rights notice: licenses/arxiv-2509.24787-CC-BY-4.0.txt.\par\smallskip

\noindent\textit{Editorial scope.} Counted unit: the Proposition whose source label is \texttt{prop:H\_Q*\_bij} in the article (source0.tex lines 677--679, source label \texttt{prop:H\_Q*\_bij}), delivered together with the article's own complete original proof of it (source0.tex lines 681--714, one \texttt{proof} environment of 34 lines). No mathematics is written, repaired or re-derived: statement and proof are the article's own text line for line. Required context retained uncounted: enumitem:sum\_condition (lines 508--520). Every definition, display and label of those lines is retained unchanged. Omitted from this package and not redistributed: 1--507, 521--676, 680--680, 715--1165. Nothing inside a retained excerpt is omitted or altered: every retained body is delivered exactly as the article's lines are written, so no omission receipt of any kind accompanies this package. Citations: the retained excerpts carry no citation command at all, so no bibliography entry of the article is transcribed and no reference is redistributed. Reference adaptation: none was needed and none was made; every reference inside the delivered unit resolves inside it, so no unresolved reference or citation is produced. Printed slips kept literal: no printed word, symbol, hypothesis or number of the counted statement or of its proof was altered. Layout and class: the article's own class is \texttt{article} with packages including inputenc, fontenc, babel, amsmath, amsfonts, amssymb, amsthm, enumitem, mathtools, booktabs, multirow, libertine, newtxmath, mathalfa; the wrapper is the lane's pinned \texttt{amsart} editorial wrapper at 10pt on A4, which restates the theorem-like environments the retained text needs and adds the article's own macro definitions that the retained text needs (\texttt{\textbackslash Z}), with extra wrapper packages (bm, bbm, dsfont, xspace, cancel, mathtools, enumitem). The delivered numbering is the unit's own. Assets: no retained span contains a figure environment, a graphics-inclusion command, a PDF-page inclusion, a table or a TikZ picture; no image file is copied into this package, the package declares no asset dependency and no image file is redistributed with it. Licence: version arXiv:2509.24787v1 of this article is distributed under the Creative Commons Attribution 4.0 International licence, as recorded in the arXiv licence metadata for this exact version and shown on the abstract page https://arxiv.org/abs/2509.24787v1. A full rights notice is delivered at licenses/arxiv-2509.24787-CC-BY-4.0.txt. Characters: every retained excerpt is pure ASCII, so the delivered engine, XeLaTeX with its default Latin Modern text font, reports no missing character.
\begin{definition}
    We define a \emph{partition tree} of degree $n$ and base-length $p$ to be a tuple $(T,v_0,f_V,f_E)$. 
    Here, $T$ is a rooted binary plane tree with $n$ leaves and root $v_0$. 
    The functions $f_V: V_I(T)\rightarrow \Z_{\geq0}$ and $f_E: E(T)\rightarrow\Z$ are integer labelings on internal vertices and the edges respectively, such that
    \begin{enumerate}[label=(P.\arabic*)]
        \item\label{enumitem:sum_condition} for all $v\in V_I(T)$ 
        \begin{align}\label{eq:fV_fE_relation}
            f_V(v)=f_E(v^\textrm{L})+f_E(v^\textrm{R}) - f_E(v^-)+1.
        \end{align}
        Heuristically, the labels on the vertices indicate the `excess' of edge labels, offset by one.
        \item\label{enumitem:root_label} $f_E(e_0)=p$ where $e_0$ is the edge incident to the root $v_0$.
    \end{enumerate}
\end{definition}

\begin{proposition}\label{prop:H_Q*_bij}
    For all $n,p$ there exists a bijection $\psi$ between $\mathcal{H}_{n,p}$ and $\mathcal{Q}^*_{n,p}$.
\end{proposition}

\begin{proof} 
    We will have to get rid of the $f_V(v)\neq0$ cases that can occur when $v$ is a bottom vertex, while keeping requirement~\ref{enumitem:sum_condition}. 
    We can do this by lowering the labels on the path that starts at $v$.  

    To be precise: 
    Let $\psi$ be the map from H-trees to well-based Q-trees by keeping the rooted tree and sending $(f_V,f_E)\mapsto (f_V',f_E')$ such that $f_V'=0$ and
    \begin{align}
        f_E'(e)=\begin{dcases}
            f_E(e) - f_V(v) &\textrm{ if } e\in P_\textrm{reg}(v)\\
            f_E(e)          &\textrm{ otherwise}
        \end{dcases}.
    \end{align}
    Note that (non-)positive edges remain (non-)positive, so that the paths are invariant under $\psi$. 
    Furthermore, strong paths remain strong, since all edge labels within a path are shifted by the same amount.
    
    It is therefore easy to see that the $\psi$-image of an H-tree is indeed a well-based Q-tree of the same base-length and degree.

    We can find the inverse $\psi^{-1}$ that maps well-based Q-trees to H-trees by, again, keeping the rooted tree and sending $(f_V,f_E)\mapsto (f_V',f_E')$ such that 
    \begin{align}
        f_V'(v)=\begin{dcases}
            -f_E(l(v)) & \textrm{ if $v$ a bottom vertex}\\  
            0 & \textrm{ otherwise}
        \end{dcases}
    \end{align}
    and
    \begin{align}
        f_E'(e)=\begin{dcases}
            f_E(e) - f_E(l(v)) &\textrm{ if } e\in P_\textrm{reg}(v)\\
            f_E(e)          &\textrm{ otherwise}
        \end{dcases},
    \end{align}
    where we remind the reader that $l(v)$ is the final edge in the path $P_\textrm{reg}(v)$.
    Also note that since the paths are strong, the paths are again preserved, and we have $f_E(e)=0$ only when $e$ is incident to a leaf.
\end{proof}

\end{document}
